\chapter{Business Functions, Business Processes, and ERP Integration}
\label{ch:business-processes}

\chapterguide
  {You only need a basic idea of what departments such as Sales, Finance, Supply Chain, and HR do.}
  {distinguish a functional area from a business function and a business process, trace cross-functional information flows, and explain why a shared ERP database improves coordination.}

\begin{sourcebasis}
This chapter follows Tutorial 1 on Business Process / Business Functions and
Chapter 1 of \emph{Concepts in Enterprise Resource Planning}. The West End
Bicycles scenario and the four functional areas are preserved from the course
materials; the explanations and diagrams reorganise those ideas into a study
sequence.
\end{sourcebasis}

\section{Start with the organisation, not the software}

An ERP system exists to support how a business works. Before learning menus and
modules, identify the work the organisation actually performs.

A \term{functional area} is a broad part of the organisation responsible for a
family of activities. A \term{business function} is one activity inside such an
area --- ``taking a sales order.'' A \term{business process} is broader: a
connected set of activities that turns an input into an output someone values.
The three are often confused because they sit at different levels.

\begin{erpfig}[Area, function, process]{The three levels. Functional areas are
vertical containers of responsibility; business functions are the individual
activities inside them; a business process runs horizontally across the
containers. The mismatch between the vertical structure and the horizontal flow
is the problem ERP exists to solve.}{fig:area-function-process}
\begin{tikzpicture}[x=1cm,y=1cm]
% --- four vertical containers ---
\foreach \x/\nm/\c in {0/Marketing \& Sales/UTPBlue, 3.9/Supply Chain/UTPTeal,
                       7.8/Accounting \& Finance/UTPGreen, 11.7/Human Resources/UTPOrange}{
  \draw[draw=\c!55,fill=\c!7,rounded corners=2mm] (\x,-0.15) rectangle (\x+3.45,4.5);
  \node[font=\scriptsize\bfseries\color{\c!75!black},align=center,text width=3.2cm]
    at (\x+1.72,4.12) {\nm};
}
% --- functions inside each ---
\foreach \x/\a/\b/\c in {0/Quotation/Sales order/Forecast,
                         3.9/Purchasing/Production/Delivery,
                         7.8/Invoicing/Receivables/Reporting,
                         11.7/Recruiting/Training/Payroll admin.}{
  \node[erpdoc,minimum width=2.9cm,minimum height=0.55cm] at (\x+1.72,3.42) {\a};
  \node[erpdoc,minimum width=2.9cm,minimum height=0.55cm] at (\x+1.72,2.72) {\b};
  \node[erpdoc,minimum width=2.9cm,minimum height=0.55cm] at (\x+1.72,2.02) {\c};
}
\node[erpcap,anchor=west,text width=2.6cm] at (-0.05,-0.75)
  {\bfseries functional\\areas};
\node[erpcap,anchor=west] at (2.9,-0.75) {each holds many \emph{business functions}};
% --- horizontal process band ---
\draw[draw=UTPNavy,fill=UTPNavy!8,rounded corners=2mm] (-0.35,0.35) rectangle (15.5,1.4);
\node[font=\scriptsize\bfseries\color{UTPNavy},anchor=west] at (-0.15,1.12)
  {Business process: \emph{order-to-cash}};
\foreach \x/\t in {1.2/Capture, 4.3/Reserve, 7.3/Ship, 10.3/Invoice, 13.3/Collect}{
  \node[erpchip] (c\t) at (\x,0.68) {\t};
}
\draw[erpflow] (cCapture) -- (cReserve); \draw[erpflow] (cReserve) -- (cShip);
\draw[erpflow] (cShip) -- (cInvoice);   \draw[erpflow] (cInvoice) -- (cCollect);
\end{tikzpicture}
\end{erpfig}

\begin{erptable}
\begin{tabularx}{\linewidth}{@{}>{\bfseries\leavevmode\color{ERPNavy}}p{0.23\linewidth}X@{}}
\toprule
\rowcolor{ERPPaleBlue!92}
Functional area & Typical business functions \\
\midrule
Marketing and Sales & Product promotion, pricing, quotations, sales orders, customer support, CRM, and sales forecasting. \\
Supply Chain Management & Purchasing, receiving, inventory, production planning, manufacturing, logistics, and delivery. \\
Accounting and Finance & Recording transactions, receivables and payables, budgeting, cash management, financial reporting, and profitability analysis. \\
Human Resources & Recruiting, hiring, training, payroll-related administration, benefits, employee records, and compliance. \\
\bottomrule
\end{tabularx}
\end{erptable}

\begin{intuition}
A department is an organisational box. A process is a journey. The customer sees
one experience; the company sees several departments.
\end{intuition}

\section{Why business processes cross functional boundaries}

Suppose West End Bicycles receives an order for 1,000 bicycles. A sales employee
can capture the request, but Sales alone cannot complete it. Six questions must
be answered, and each belongs to a different desk.

\begin{erpfig}[Order questions by owner]{One customer order, six questions, four
owners. No single functional area can answer all of them, which is what makes
this a cross-functional process rather than a departmental
task.}{fig:cross-functional-order}
\begin{tikzpicture}[x=1cm,y=1cm]
\node[erphead,anchor=west] at (-0.2,0.75) {Order for 1{,}000 bicycles --- who answers what?};
\foreach \i/\q/\own/\c in {%
   0/{Is the order captured with the right product, quantity, and price?}/{Marketing \& Sales}/UTPBlue,
   1/{Is the customer allowed to buy on the requested credit terms?}/{Accounting \& Finance}/UTPGreen,
   2/{Are 1{,}000 bicycles already in stock?}/{Supply Chain}/UTPTeal,
   3/{If not, can production make them on time from available parts?}/{Supply Chain}/UTPTeal,
   4/{When can the order be delivered?}/{Supply Chain}/UTPTeal,
   5/{What is invoiced, and has the customer paid?}/{Accounting \& Finance}/UTPGreen}{
  \pgfmathsetmacro{\y}{-\i*0.92}
  \node[erpdoc,anchor=west,text width=8.5cm,minimum width=9.0cm,minimum height=0.74cm,
        align=left,draw=\c!60] (q\i) at (0,\y) {\q};
  \node[erpstep,anchor=west,minimum width=3.7cm,minimum height=0.74cm,
        draw=\c,fill=\c!10] at (9.6,\y) {\own};
  \draw[erplink] (9.05,\y) -- (9.6,\y);
}
\foreach \i in {0,...,4}{
  \pgfmathsetmacro{\y}{-\i*0.92-0.37}
  \draw[erpflow] (0.55,\y) -- (0.55,\y-0.18);
}
\draw[draw=UTPMuted!35] (0,-5.05) -- (13.3,-5.05);
\node[erpstep,anchor=west,minimum width=3.7cm,minimum height=0.74cm,
      draw=UTPOrange,fill=UTPPaleOrange] (hr) at (9.6,-5.55) {Human Resources};
\node[erpcap,anchor=west,text width=8.5cm,align=left] at (0,-5.55)
  {\itshape staffs every step above, owns none of them};
\draw[erplink] (9.05,-5.55) -- (hr);
\end{tikzpicture}
\end{erpfig}

HR may not touch an individual order, yet it supports the process by ensuring
the company has trained salespeople, buyers, planners, accountants, and
warehouse staff.

\section{The information each area needs from the others}

A process fails when information moves slowly or inconsistently, even if every
department works hard. Each area is simultaneously a producer and a consumer of
the others' data.

\begin{erpfig}[Information exchange]{What each functional area produces, and
what it must receive from the others. Every labelled arrow is a dependency that
a silo turns into a delay.}{fig:info-exchange}
\begin{tikzpicture}[x=1cm,y=1cm]
\node[erpstep,minimum width=3.6cm,minimum height=1.15cm] (ms)  at (0,3.1)
  {Marketing \& Sales\\[-0.2em]{\tiny\mdseries produces orders, prices, demand}};
\node[erpstep,minimum width=3.6cm,minimum height=1.15cm,draw=UTPTeal,fill=UTPPaleTeal] (scm) at (9.6,3.1)
  {Supply Chain\\[-0.2em]{\tiny\mdseries produces stock, output, shipments}};
\node[erpstep,minimum width=3.6cm,minimum height=1.15cm,draw=UTPGreen,fill=UTPPaleTeal] (af) at (9.6,0)
  {Accounting \& Finance\\[-0.2em]{\tiny\mdseries produces invoices, statements}};
\node[erpstep,minimum width=3.6cm,minimum height=1.15cm,draw=UTPOrange,fill=UTPPaleOrange] (hr) at (0,0)
  {Human Resources\\[-0.2em]{\tiny\mdseries produces staff, records, skills}};
% top: sales -> scm and back
\draw[erpflow] ([yshift=3mm]ms.east) -- ([yshift=3mm]scm.west)
  node[midway,above,erpnote]{forecasts, confirmed orders};
\draw[erpback] ([yshift=-3mm]scm.west) -- ([yshift=-3mm]ms.east)
  node[midway,below,erpnote]{availability, promised dates};
% right: scm -> finance
\draw[erpflow] (scm.south) -- (af.north)
  node[midway,right,erpnote,text width=2.7cm]{receipts, production, shipments};
% bottom: finance -> sales
\draw[erpback] ([yshift=3mm]af.west) -- ([yshift=3mm]hr.east)
  node[midway,above,erpnote]{payroll and expense data};
\draw[erpflow] (af.north west) to[out=150,in=-30]
  node[midway,fill=white,inner sep=1pt,erpnote,text width=2.4cm,align=center]
  {credit status, payments} (ms.south east);
% left: hr -> sales
\draw[erpflow] (hr.north) -- (ms.south)
  node[midway,left,erpnote,text width=2.2cm]{trained staff, skills};
\draw[erpback] ([yshift=-3mm]hr.east) -- ([yshift=-3mm]af.west)
  node[midway,below,erpnote]{headcount, staffing needs};
\end{tikzpicture}
\end{erpfig}

\begin{keyidea}
The quality of a decision is limited by the quality, timeliness, and consistency
of the data available when the decision is made. ERP integration is valuable
because it removes the delay and disagreement created by several disconnected
copies of the same business fact.
\end{keyidea}

\section{Silos versus an integrated information system}

A \term{silo} is a departmental system that keeps its own data and does not
automatically stay synchronised with the others. Silos are not inherently bad.
The problem appears when a cross-functional process depends on a fact that now
exists in several different versions.

\begin{erpfig}[Silos versus shared data]{The same three departments, twice. On
the left each holds its own copy and the copies drift apart between transfers;
on the right one record serves all three, so there is nothing to
reconcile.}{fig:silo-vs-erp}
\resizebox{0.99\linewidth}{!}{%
\begin{tikzpicture}[x=1cm,y=1cm]
% ---- left: silos ----
\node[erphead,color=UTPRed] at (3.6,3.55) {Disconnected silos};
\node[erpbad,minimum width=2.6cm] (s1) at (0.0,2.3) {Sales DB};
\node[erpbad,minimum width=2.6cm] (s2) at (3.6,2.3) {Warehouse DB};
\node[erpbad,minimum width=2.6cm] (s3) at (7.2,2.3) {Finance DB};
\draw[erpslow] (s1) -- node[erpnote,above]{fax / batch} (s2);
\draw[erpslow] (s2) -- node[erpnote,above]{re-keyed} (s3);
\node[erpdoc,minimum width=2.4cm,minimum height=0.6cm,draw=UTPRed] (c1) at (0.0,0.85) {qty $=$ 40};
\node[erpdoc,minimum width=2.4cm,minimum height=0.6cm,draw=UTPRed] (c2) at (3.6,0.85) {qty $=$ 38};
\node[erpdoc,minimum width=2.4cm,minimum height=0.6cm,draw=UTPRed] (c3) at (7.2,0.85) {qty $=$ 40};
\foreach \a/\b in {s1/c1,s2/c2,s3/c3}{\draw[erplink] (\a)--(\b);}
\node[erpchipbad] at (3.6,0.05) {three answers to one question};
% ---- right: integrated ----
\begin{scope}[xshift=10.6cm]
\node[erphead,color=UTPTeal] at (3.1,3.55) {Integrated ERP};
\node[erpgood,minimum width=2.6cm] (i1) at (0.0,2.3) {Sales};
\node[erpgood,minimum width=2.6cm] (i2) at (3.1,2.3) {Inventory};
\node[erpgood,minimum width=2.6cm] (i3) at (6.2,2.3) {Finance};
\node[erpdb,minimum width=3.0cm,minimum height=1.2cm] (db) at (3.1,0.7) {Shared record\\qty $=$ 40};
\foreach \n in {i1,i2,i3}{\draw[erpflow] (\n) -- (db);}
\node[erpchip] at (3.1,-0.42) {one answer, read three ways};
\end{scope}
\end{tikzpicture}}
\end{erpfig}

The structural cost of silos is easiest to see at the keyboard: the same fact
gets typed once per system, and each retype is an opportunity to disagree.

\begin{erpfig}[Cost of re-entry]{Where errors enter. Every manual hand-off is a
chance for the copies to diverge; the integrated path has one entry point and
no reconciliation step.}{fig:duplicate-entry}
\resizebox{\linewidth}{!}{%
\begin{tikzpicture}[x=1cm,y=1cm]
\node[erpcap,anchor=east,text width=1.9cm] at (-0.15,1.5) {\bfseries Siloed};
\node[erpbad,minimum width=2.1cm,minimum height=0.72cm] (a1) at (1.2,1.5) {type};
\node[erpbad,minimum width=2.1cm,minimum height=0.72cm] (a2) at (4.3,1.5) {re-type};
\node[erpbad,minimum width=2.1cm,minimum height=0.72cm] (a3) at (7.4,1.5) {re-type};
\node[erpwarn,minimum width=2.4cm,minimum height=0.72cm] (a4) at (10.7,1.5) {reconcile};
\draw[erpslow] (a1)--(a2); \draw[erpslow] (a2)--(a3); \draw[erpslow] (a3)--(a4);
\node[erpchipbad,anchor=west] at (12.15,1.5) {3 entries, 3 error points};
\node[erpcap,anchor=east,text width=1.9cm] at (-0.15,0.0) {\bfseries Integrated};
\node[erpgood,minimum width=2.1cm,minimum height=0.72cm] (b1) at (1.2,0.0) {type once};
\node[erpdb,minimum width=2.4cm,minimum height=0.9cm] (b2) at (4.6,0.0) {shared};
\node[erpgood,minimum width=2.1cm,minimum height=0.72cm] (b3) at (7.9,0.0) {read};
\node[erpgood,minimum width=2.1cm,minimum height=0.72cm] (b4) at (10.7,0.0) {read};
\draw[erpflow] (b1)--(b2); \draw[erpflow] (b2)--(b3); \draw[erpflow] (b3)--(b4);
\node[erpchip,anchor=west] at (12.2,0.0) {1 entry, 0 error points};
\end{tikzpicture}}
\end{erpfig}

With a shared database one business event updates several views at once. When a
delivery is validated, Inventory sees stock leave, Sales sees fulfilment
progress, and Finance gains the evidence needed to bill.

\section{West End Bicycles as the running example}

The tutorial scenario describes West End Bicycles as a manufacturer of road,
mountain, and city bicycles plus accessories. It is useful precisely because it
requires all four functional areas at once.

\begin{erpfig}[Two connected cycles]{The whole course in one picture. The buying
cycle fills the warehouse, the selling cycle empties it, and manufacturing is
the hinge between them. Every later chapter deepens one segment of this
loop.}{fig:two-cycles}
\begin{tikzpicture}[x=1cm,y=1cm]
\node[draw=UTPBlue!35,fill=UTPPaleBlue!45,rounded corners=2.5mm,
      minimum width=9.6cm,minimum height=1.8cm] at (4.2,2.35) {};
\node[erpcap,anchor=west,color=UTPBlue] at (-0.55,3.05) {\bfseries Procure-to-pay};
\node[erpstep,minimum width=2.05cm] (n1) at (0.75,2.25) {Need};
\node[erpstep,minimum width=2.05cm] (n2) at (3.05,2.25) {Purchase order};
\node[erpstep,minimum width=2.05cm] (n3) at (5.35,2.25) {Receipt};
\node[erpstep,minimum width=2.05cm] (n4) at (7.85,2.25) {Pay vendor};
\draw[erpflow] (n1)--(n2); \draw[erpflow] (n2)--(n3); \draw[erpflow] (n3)--(n4);

\node[erpwarn,minimum width=3.4cm,minimum height=1.05cm] (mfg) at (12.1,1.05)
  {Manufacturing\\[-0.15em]{\tiny\mdseries components $\rightarrow$ Road Bikes}};

\node[draw=UTPGreen!35,fill=UTPPaleTeal!55,rounded corners=2.5mm,
      minimum width=9.6cm,minimum height=1.8cm] at (4.2,-0.25) {};
\node[erpcap,anchor=west,color=UTPGreen] at (-0.55,-1.15) {\bfseries Order-to-cash};
\node[erpgood,minimum width=2.05cm] (m1) at (0.75,-0.35) {Opportunity};
\node[erpgood,minimum width=2.05cm] (m2) at (3.05,-0.35) {Sales order};
\node[erpgood,minimum width=2.05cm] (m3) at (5.35,-0.35) {Delivery};
\node[erpgood,minimum width=2.05cm] (m4) at (7.85,-0.35) {Invoice / payment};
\draw[erpflow] (m1)--(m2); \draw[erpflow] (m2)--(m3); \draw[erpflow] (m3)--(m4);

\draw[erpflow] (n4) to[bend left=16] node[erpnote,above right=0pt and 1pt]{materials in} (mfg.north west);
\draw[erpflow] (mfg.south west) to[bend left=16] node[erpnote,below right=0pt and 1pt]{bikes out} (m4.north east);
\end{tikzpicture}
\end{erpfig}

\section{What Odoo adds to this picture}

The supplied Odoo 19 Community guide implements the same concepts with eight
applications. Reading the app grid as an architecture rather than a menu is the
fastest way to see where each chapter lands.

\begin{erpfig}[Odoo apps by functional area]{The eight course applications
grouped by the functional area each serves. Contacts sits underneath all of
them because customer and vendor identity is shared master
data.}{fig:odoo-app-map}
\begin{tikzpicture}[x=1cm,y=1cm]
\foreach \x/\nm/\c in {0/Marketing \& Sales/UTPBlue, 5.0/Supply Chain/UTPTeal,
                       10.0/Accounting \& Finance/UTPGreen}{
  \draw[draw=\c!55,fill=\c!7,rounded corners=2mm] (\x,0.15) rectangle (\x+4.5,2.75);
  \node[font=\scriptsize\bfseries\color{\c!75!black}] at (\x+2.25,2.43) {\nm};
}
\draw[draw=UTPOrange!55,fill=UTPOrange!7,rounded corners=2mm] (10.0,-2.4) rectangle (14.5,-0.5);
\node[font=\scriptsize\bfseries\color{UTPOrange!75!black}] at (12.25,-0.85) {Human Resources};
\node[erpstep,minimum width=1.85cm] at (1.15,1.65) {CRM};
\node[erpstep,minimum width=1.85cm] at (3.35,1.65) {Sales};
\node[erpstep,minimum width=1.85cm,draw=UTPTeal,fill=UTPPaleTeal] at (6.15,1.65) {Purchase};
\node[erpstep,minimum width=1.85cm,draw=UTPTeal,fill=UTPPaleTeal] at (8.35,1.65) {Inventory};
\node[erpstep,minimum width=4.0cm,draw=UTPTeal,fill=UTPPaleTeal] at (7.25,0.72) {Manufacturing};
\node[erpstep,minimum width=3.9cm,draw=UTPGreen,fill=UTPPaleTeal] at (12.25,1.65) {Invoicing};
\node[erpcap,text width=3.9cm] at (12.25,0.75) {invoices and payments in the Community workflow};
\node[erpstep,minimum width=3.9cm,draw=UTPOrange,fill=UTPPaleOrange] at (12.25,-1.55) {Employees};
\node[erpstore,minimum width=8.4cm,minimum height=0.95cm] (ct) at (4.6,-1.55)
  {Contacts --- shared customer and vendor identity};
\draw[erplink] (ct) -- (1.15,0.15); \draw[erplink] (ct) -- (7.25,0.15);
\draw[erplink] (ct) -- (10.6,0.15);
\end{tikzpicture}
\end{erpfig}

\begin{odoobox}
Three ideas recur throughout the practicals:
\begin{itemize}
  \item Records are linked. A purchase order leads to a receipt; a sales order leads to a delivery and an invoice.
  \item Documents move through states such as draft, confirmed/posted, and done/paid.
  \item Smart buttons expose those relationships, making the information flow visible instead of requiring manual reconciliation.
\end{itemize}
\end{odoobox}

\begin{pitfall}
Do not define ERP as ``software that combines departments.'' The point is
process integration: one business event should be visible consistently to the
functions that need it, when they need it.
\end{pitfall}

\begin{quickcheck}
\begin{enumerate}
  \item Explain the difference between a functional area, a business function, and a business process using West End Bicycles.
  \item Why can a weekly printed credit report cause a sales-order problem even if the Accounting data was correct when it was printed?
  \item Name three pieces of information that must move between Sales and Supply Chain when fulfilling a bicycle order.
\end{enumerate}
\end{quickcheck}

\begin{chapterrecap}
\begin{itemize}
  \item Functional areas organise responsibility; business processes connect activities across them (\figref{fig:area-function-process}).
  \item The four areas depend on one another's data in both directions (\figref{fig:info-exchange}).
  \item Disconnected systems create duplicate entry, delay, inconsistency, and poor visibility (\figref{fig:silo-vs-erp}).
  \item ERP lets one business event serve several departments without re-entry.
  \item West End Bicycles runs two linked cycles --- procure-to-pay and order-to-cash --- joined by manufacturing (\figref{fig:two-cycles}).
\end{itemize}
\end{chapterrecap}
